\chapter{Theoretical Background}

\section{Basics of Cryptography} 

Cryptography, derived from Greek words meaning "hidden writing", is the study
and practice of mathematical methods for secure communication over insecure
channels. In modern systems, it is no longer limited to confidentiality alone,
but also addresses integrity, authentication, and non-repudiation. More
broadly, modern cryptography is concerned not only with secrecy, but also with
detecting modification and verifying
authenticity~\cite{diffie1976newDirections,ibm2023cryptography}.

The need for secret communication predates modern computing by thousands of
years. One early example is the use of unusual hieroglyphs carved into the wall
of a tomb from the Old Kingdom of Egypt, often cited as one of the earliest
known uses of cryptography~\cite{schneider2024ibmhistory}.

Before computers, cryptography relied on classical ciphers that operated on
letters rather than bits. One of the most popular examples is the Caesar
cipher, which shifts each letter by a fixed number of positions in the
alphabet. The cipher is obviously easy to break by trying all possibilities
because number of possible shifts is really small. A more advanced classical
cipher is the Vigen\`ere cipher, which uses a keyword to vary the shift applied
to each letter, but these ciphers are still vulnerable to frequency and
statistical analysis of the text. Classical ciphers were historically important, but their
security was too weak for modern communication
systems~\cite{aumasson2017serious}.

Later, more sophisticated systems were developed, including mechanical rotor
machines such as the Enigma machine used for scrambling war messages during
World War II. These devices used changing and more sophisticated substitutions,
which made them harder to break than earlier hand
ciphers~\cite{schneider2024ibmhistory}.

Their eventual defeat showed that security could not rely only on increasingly
complicated substitutions. From the 1970s onward, cryptography began to move
toward mathematically defined constructions whose security depended on the
difficulty of underlying computational problems~\cite{diffie1976newDirections}. This shift marks the beginning
of modern cryptography and motivates the discussion in the next
section.

\section{P, NP, and Computational Hardness}

The security of many modern cryptosystems depends on the assumption that a
related computational problem is hard to solve efficiently. Computational
complexity theory provides a way to describe computational difficulty. The
following concepts introduce the basic language used to discuss this.

To describe computational difficulty more precisely, problems are organized
into complexity classes. They reflect how the running time of an algorithm
grows with the size of the input. In complexity theory, $\mathcal{P}$ and
$\mathcal{NP}$ are classes of decision problems, which are computational
problems that ask whether a given input satisfies a specified condition. Each
valid input is answered by either yes or
no.~\cite[pp.~18--19]{garey1979computers} 

The class $\mathcal{P}$ contains those decision problems that can be solved in
polynomial time with respect to the input size, and they are commonly regarded
as efficiently solvable. Here, efficient means solvable in polynomial time with
respect to the input size~\cite[p.~27]{garey1979computers}.

The class $\mathcal{NP}$, which stands for nondeterministic polynomial time,
contains decision problems for which a proposed solution can be checked in
polynomial time. This does not mean that the solution can also be found
efficiently. It only means that, if someone gives us a candidate solution, we
can verify it efficiently.~\cite[pp.~28--29]{garey1979computers}

A major open problem in complexity theory is whether $\mathcal{P} =
\mathcal{NP}$. If this were true, every problem whose solution can be verified
in polynomial time could also be solved in polynomial time. However, no proof
is currently known for either $\mathcal{P} = \mathcal{NP}$ or $\mathcal{P} \neq
\mathcal{NP}$. Most researchers nevertheless believe that $\mathcal{P} \neq
\mathcal{NP}$.~\cite[ch.~9]{aumasson2017serious}

To compare the difficulty of decision problems, complexity theory uses
polynomial-time reductions. Informally, a polynomial-time reduction is an
efficient transformation from one problem to another that preserves the answer
to each instance. If one problem can be reduced to another in polynomial time,
then an efficient algorithm for the second problem would also give an efficient
algorithm for the first. Reductions therefore provide a standard way to compare
computational difficulty between problems.

A decision problem is called $\mathcal{NP}$-complete if it belongs to
$\mathcal{NP}$ and every problem in $\mathcal{NP}$ can be reduced to it in
polynomial time. If any $\mathcal{NP}$-complete problem were solvable in
polynomial time, then every problem in $\mathcal{NP}$ would also be solvable in
polynomial time.

A related notion is $\mathcal{NP}$-hard. Informally, a problem is
$\mathcal{NP}$-hard if it is at least as difficult as the
$\mathcal{NP}$-complete problems under polynomial reductions. The
$\mathcal{NP}$-complete problems form the part of this broader class that
also belongs to $\mathcal{NP}$.~\cite[pp.~34--35, 37]{garey1979computers}

Modern cryptography often relies on mathematical problems for which no
efficient algorithm is known. This is especially important in public-key
cryptography, where security depends on assumptions about computational
hardness rather than on keeping the method secret. Different cryptosystems rely
on different computational problems, such as integer factorization, discrete
logarithms, or subset-sum.

One problem especially relevant for this thesis is the subset-sum problem,
whose decision version is
$\mathcal{NP}$-complete.~\cite[p.~223]{garey1979computers} The hardness of
this problem was one of the reasons why knapsack-type constructions were
originally considered promising for cryptographic use. The subset-sum problem
and the Merkle--Hellman construction are introduced later in
Sections~\ref{sec:subset-sum-problem} and~\ref{sec:merkle-hellman}.

However, this connection must be interpreted carefully.
$\mathcal{NP}$-completeness refers to the worst-case hardness of the general
problem, whereas a cryptosystem uses specific structured instances. Diffie and
Hellman~\cite[p.~654]{diffie1976newDirections} already noted that complexity
theory gives no indication of the difficulty of a particular knapsack
instance. This distinction will be
important in the later discussion of knapsack cryptosystems.

\section{Symmetric and Asymmetric Cryptography}

Modern cryptography is usually divided into two main branches. These are
symmetric cryptography, also called private-key cryptography, and asymmetric
cryptography, also called public-key cryptography. Both are used to protect
information, but they diverge mainly in how cryptographic keys are used. This
leads to different practical advantages, limitations, and security assumptions.

In symmetric cryptography, the same secret key is used for both encryption and
decryption. All communicating parties must use the same key and keep it secret.
This means that the key must be established and exchanged before communication.
Because of this limitation, symmetric systems are often not practical when two
parties have never communicated before or do not have a secure channel for
exchanging the key. Despite this, it is still widely used in practice and
especially useful for encrypting large amounts of data efficiently. Modern
symmetric encryption commonly uses block ciphers or stream ciphers, with AES
being a well-known block cipher example.~\cite{ibm2023cryptography}

This problem motivated the development of a different approach, in which the
parties do not need to share a secret beforehand. This idea led to public-key
cryptography. It works with a pair of keys, a public key that can be shared
with anyone and a private key that must remain
secret~\cite[pp.~644--645]{diffie1976newDirections}. Messages encrypted with
the public key can only be decrypted with the corresponding private key. This
removes the need to exchange a secret encryption key through a secure channel
beforehand, although the public key itself still has to be authenticated in
practice. Diffie and Hellman~\cite[pp.~644, 648--649]{diffie1976newDirections}
also described a key exchange protocol in which two users
communicating over a public channel can obtain a shared secret key, which can
then be used in a conventional symmetric cryptosystem. This shows that
symmetric and asymmetric cryptography do not replace one another, and can be
used together.

The security of asymmetric systems is based on the idea of one-way and trapdoor
functions. Diffie and Hellman~\cite[pp.~650, 652--654]{diffie1976newDirections}
described these as functions that are easy to compute, but difficult to invert
without additional secret information, and suggested the knapsack problem as
one possible foundation for such systems. In the cryptosystem studied in this
thesis, the underlying problem is subset-sum, also known as the knapsack
problem. The following sections introduce this problem and explain how it can
be used to construct a public-key encryption scheme.

\section{The Subset-Sum Problem}
\label{sec:subset-sum-problem}

The subset-sum problem, often also referred to as the knapsack problem, is
defined as follows. Given a sequence of positive integers
\[
W = (w_1, w_2, \dots, w_n)
\]
and a target value $s \in \mathbb{N}$, the search version asks to find a binary
vector
\[
x = (x_1, x_2, \dots, x_n) \in \{0,1\}^n
\]
such that
\[
\sum_{i=1}^{n} x_i w_i = s.
\]
Here, each variable $x_i$ indicates whether the weight $w_i$ is included in the
subset or not. The related decision problem asks whether such a vector exists.
This decision version is
$\mathcal{NP}$-complete~\cite[p.~223]{garey1979computers}.

The distinction between the decision and search versions is important for
cryptography. In encryption, the attacker is not only interested in if a
solution exists. The attacker needs to recover the actual vector $x$, because
this vector represents the encrypted message.~\cite[pp.~75--76]{OdlyzkoRiseFall}

\subsection*{Example}

Consider the sequence
\[
W = (3,\, 8,\, 12,\, 15,\, 19)
\]
and the target value
\[
s = 27.
\]
To check if some subset sums to $27$, we can examine different combinations
of the weights. In this case, there are two valid solutions. The first one is
\[
8 + 19 = 27,
\]
which corresponds to the vector
\[
x = (0,\,1,\,0,\,0,\,1).
\]
The second one is
\[
12 + 15 = 27,
\]
which corresponds to
\[
x = (0,\,0,\,1,\,1,\,0).
\]
This example is intentionally small, so the definition is easy to understand.
It does not mean that the general problem is easy. For a sequence of length
$n$, there are $2^n$ possible subsets, and the decision version of subset-sum
is $\mathcal{NP}$-complete, as stated above.

\subsection*{Simple solving methods}

A simple way to solve the search version of subset-sum is brute-force search.
This method checks possible binary vectors $x\in\{0,1\}^n$ one by one and
computes whether
\[
    \sum_{i=1}^{n} x_i w_i = s.
\]
In the worst case, it may have to check all $2^n$ vectors, so its running time
grows exponentially with the length of the instance. This makes it useful only
for small examples and experiments.

A better method is the meet-in-the-middle algorithm. It splits the weights into
two parts. First, it computes all possible subset sums of the first half. For
the second half, it computes all values obtained by subtracting each possible
subset sum of the second half from $s$. Creating these two lists already takes
exponential time, because each half has roughly $2^{n/2}$ possible subsets.
After the two lists are sorted, they are scanned for a common value. If the
same value appears in both lists, then the subset from the first half has the
same sum as $s$ minus the subset sum from the second half. Together, these two
subsets therefore sum to $s$. The total running time is roughly $n2^{n/2}$
operations. The disadvantage is that the partial sums have to be stored, so the
method also needs $2^{n/2}$ storage space~\cite[pp.~75--76]{OdlyzkoRiseFall}.

There are more advanced subset-sum algorithms and time-memory trade-offs, but
they are not the focus of this thesis. The implementation part uses brute-force
search and meet-in-the-middle search as simple comparison methods for message
recovery on small generated instances.

\subsection*{Superincreasing sequences}

Superincreasing sequences form a special easy case of subset-sum. This is
important for the Merkle--Hellman cryptosystem because the legitimate receiver
needs a private subset-sum instance that can be solved quickly. A sequence $W$
is called \emph{superincreasing} if every element is strictly larger than the
sum of all previous elements:
\[
    w_i > \sum_{j=1}^{i-1} w_j \qquad \text{for } i \ge 2.
\]

For superincreasing sequences, subset-sum can be solved by the greedy
algorithm. Starting from the largest weight, the algorithm includes the current
weight if it is not larger than the remaining target value, and otherwise skips
it. This works because all smaller weights together have a smaller sum than the
current weight. If the current weight is needed, the smaller weights cannot
replace it.

For example, let
\[
    W=(2,\,5,\,11,\,23,\,42)
\]
and let the target be $s=28$. The greedy algorithm starts from the largest
weight. It skips $42$, because it is larger than the remaining target value.
It then includes $23$, so the remaining value becomes $5$. Next it skips $11$,
includes $5$, and ends with remaining value $0$. Thus
\[
    28 = 23 + 5,
\]
which corresponds to the vector
\[
    x=(0,\,1,\,0,\,1,\,0).
\]
This example is again intentionally small, so the steps of the greedy algorithm
are easy to follow.

A superincreasing sequence ensures that only one solution can exist for a given
target, because the largest selected weight cannot be replaced by any
combination of smaller weights. This is different from general subset-sum
instances, where several subsets may have the same target sum. The fact that
superincreasing instances are easy to solve and have unique solutions is why
they are useful in the Merkle--Hellman cryptosystem. The private key contains a
superincreasing instance, while the public key is intended to look like a
general subset-sum instance~\cite[pp.~525--526]{MerkleHellman1978}\cite[pp.~77--79]{OdlyzkoRiseFall}.
This idea is explained in more detail in the next section.

\section{The Merkle--Hellman Knapsack Cryptosystem}
\label{sec:merkle-hellman}

The Merkle--Hellman cryptosystem was one of the earliest public-key encryption
schemes. It is based on the subset-sum problem.~\cite{MerkleHellman1978} Its
design uses a superincreasing sequence as part of the private key. This easy
instance is then transformed into a public sequence that is intended to look
like a hard general subset-sum instance. Although the idea was influential,
later cryptanalysis showed that the basic construction is
insecure.~\cite{Shamir1982}

The main parameters of the cryptosystem are the block length $n$, the private superincreasing
sequence $W=(w_1,\dots,w_n)$, the modulus $m$, the multiplier $r$, and
optionally a secret permutation $\pi$. The value $n$ determines both the number
of public weights and the number of message bits encrypted in one block.

\subsection{Key Generation}

Key generation starts from an easy private subset-sum instance and transforms
it into a public sequence by a modular transformation intended to make it look
hard. First, a superincreasing sequence $W=(w_1,\dots,w_n)$ is chosen as a
private sequence. Then a modulus $m$ larger than $\sum_{i=1}^n w_i$ and a
multiplier $r$ are selected so that $\gcd(r,m)=1$, where $\gcd$ means the
greatest common divisor. This means that $r$ and $m$ are coprime and guarantees
that $r$ has a multiplicative inverse modulo $m$~\cite[pp.~212--213]{stanoyevitch2011}.

The inverse can be found in polynomial time by the extended Euclidean
algorithm~\cite[pp.~216--217, p.~492]{stanoyevitch2011}. This algorithm finds
integers $a$ and $b$ such that
\[
    ar + bm = \gcd(r,m).
\]
When $\gcd(r,m)=1$, this gives
\[
    ar + bm = 1.
\]
Modulo $m$, the term $bm$ disappears because it is a multiple of $m$. Therefore
\[
    ar \equiv 1 \pmod m.
\]
This is exactly the condition for $a$ to be the inverse of $r$ modulo $m$, or
equivalently $r^{-1}\equiv a \pmod m$~\cite[pp.~216--217]{stanoyevitch2011}.

In the original description, the order of the transformed sequence may also be
permuted to further obscure its structure.~\cite[p.~526]{MerkleHellman1978} The
public key is the final transformed sequence, while the private key consists of
the original superincreasing sequence together with the secret modulus $m$,
multiplier $r$, and, if scrambling is used, the corresponding permutation.

\subsection{Encryption}

After key generation, the sender uses the public sequence
\[
B=(b_1,b_2,\dots,b_n)
\]
to encrypt an $n$-bit message
\[
x=(x_1,x_2,\dots,x_n), \quad x_i \in \{0,1\}.
\]
The length $n$ of the public sequence determines the block size of the scheme,
so one application of the encryption algorithm encrypts one $n$-bit block.
Longer messages must therefore be divided into blocks of this length.
The ciphertext is computed as
\[
C=\sum_{i=1}^{n} x_i b_i.
\]
In other words, only those public weights corresponding to message bits equal
to $1$ are added. Encryption is therefore straightforward and efficient, since
it requires only additions over public weights.~\cite{MerkleHellman1978}

\subsection{Decryption}

To decrypt a ciphertext $C$, the receiver first computes
\[
S' \equiv r^{-1} C \pmod m,
\]
where $r^{-1}$ is the modular inverse of $r$ modulo $m$. This multiplication
removes the multiplier from the public weights and maps the ciphertext back to
a subset-sum instance over the private superincreasing sequence.

Since $b_i \equiv r w_i \pmod m$ and
\[
    C=\sum_{i=1}^{n} x_i b_i,
\]
we get
\[
    r^{-1}C \equiv r^{-1}\sum_{i=1}^{n}x_i b_i
    \equiv \sum_{i=1}^{n}x_i w_i \pmod m.
\]

The modulus was chosen so that $m > \sum_{i=1}^{n} w_i$. Therefore every
possible private subset sum is smaller than $m$, and the value obtained after
the modular reduction is the actual private subset sum. The receiver can then
recover the message bits by the greedy algorithm. If a permutation is used, the
bit positions are restored according to the secret permutation.

\subsection*{Example}

For decryption, the inverse of $r=9$ modulo $47$ is needed. To find it, we use
the extended Euclidean algorithm. First, the Euclidean algorithm repeatedly
divides the larger number by the smaller one and keeps the remainder. Dividing
$47$ by $9$ gives quotient $5$ and remainder $2$, so
\[
    47 = 5\cdot9 + 2.
\]
Then dividing $9$ by $2$ gives quotient $4$ and remainder $1$, so
\[
    9 = 4\cdot2 + 1.
\]
The last non-zero remainder is $1$, so $\gcd(9,47)=1$.

The extended part now rewrites this remainder as a combination of $9$ and
$47$. From the second division,
\[
    1 = 9 - 4\cdot2.
\]
From the first division,
\[
    2 = 47 - 5\cdot9.
\]
Substituting this value of $2$ gives
\[
    1 = 9 - 4(47 - 5\cdot9) = 21\cdot9 - 4\cdot47.
\]
Modulo $47$, the term $-4\cdot47$ disappears. Therefore
\[
    21\cdot9 \equiv 1 \pmod{47},
\]
so $r^{-1}\equiv 21 \pmod{47}$.

Consider a private superincreasing sequence
\[
    W=(2,\,5,\,11,\,23),
\]
with $m=47$ and $r=9$, where $\gcd(r,m)=1$ and $m > 2+5+11+23 = 41$.
The public key is computed as
\[
    \begin{aligned}
        9\cdot 2 &\equiv 18 \pmod{47},\\
        9\cdot 5 &\equiv 45 \pmod{47},\\
        9\cdot 11 &\equiv 5 \pmod{47},\\
        9\cdot 23 &\equiv 19 \pmod{47}.
    \end{aligned}
\]
Thus
\[
    B=(18,\,45,\,5,\,19).
\]
For message $x=(1,0,1,0)$, encryption gives
\[
    C=18+5=23.
\]
This ciphertext can be sent over the public channel.

For decryption, the inverse of $r=9$ modulo $47$ is needed. The Euclidean
divisions are
\[
    47 = 5\cdot 9 + 2,\qquad 9 = 4\cdot 2 + 1.
\]
The second division gives $1=9-4\cdot2$. From the first division,
$2=47-5\cdot9$. Substituting this value of $2$ gives
\[
    1 = 9 - 4(47 - 5\cdot 9) = 21\cdot 9 - 4\cdot 47.
\]
Modulo $47$, the term $-4\cdot47$ disappears. Therefore
\[
    21\cdot 9 \equiv 1 \pmod{47},
\]
so $r^{-1}\equiv 21 \pmod{47}$. Then
\[
    S' \equiv r^{-1}C \equiv 21\cdot 23 \equiv 13 \pmod{47}.
\]
This is the private subset sum. Greedy recovery on $W$ gives
\[
    13=11+2,
\]
so the recovered bits are $x=(1,0,1,0)$.

This is only a toy example chosen to show the arithmetic by hand. Real
cryptographic instances and the experiments in the implementation use much
larger block lengths.

The basic construction does not hide the superincreasing structure well enough,
and later cryptanalysis showed that an equivalent private key can be recovered
in polynomial time.~\cite{Shamir1982} The attacks behind this result are
discussed in the next chapter.
